\documentclass[11pt,a4paper]{article}
\usepackage{fontspec,amsmath,amssymb,amsthm,microtype}
\setmainfont{lmroman10-regular.otf}[BoldFont=lmroman10-bold.otf,ItalicFont=lmroman10-italic.otf,BoldItalicFont=lmroman10-bolditalic.otf]
\usepackage{xcolor}
\usepackage[most]{tcolorbox}
\usepackage{mathrsfs}
\providecommand{\lt}{<}
\providecommand{\gt}{>}
\setmonofont{DejaVuSansMono.ttf}[BoldFont=DejaVuSansMono-Bold.ttf,ItalicFont=DejaVuSansMono-Oblique.ttf,BoldItalicFont=DejaVuSansMono-BoldOblique.ttf,Scale=0.86]
\newfontfamily\leanfallback{STIXTwoMath-Regular.otf}[Scale=0.86]
\definecolor{leanrule}{gray}{0.62}
\definecolor{leanback}{gray}{0.962}
\definecolor{leantitle}{gray}{0.22}
\definecolor{leanmuted}{gray}{0.38}
\newtcolorbox{leanbox}[1]{enhanced,breakable,colback=leanback,colframe=leanrule,
  boxrule=0.4pt,arc=0.6mm,left=2.4mm,right=2.4mm,top=1.8mm,bottom=1.8mm,
  before skip=0.8em,after skip=1em,fonttitle=\small\bfseries,coltitle=leantitle,
  colbacktitle=leanback,titlerule=0.3pt,title={#1}}
\newlength{\leancharwidth}
\newenvironment{leancode}{\par\smallskip\begingroup\ttfamily\footnotesize\raggedright
  \settowidth{\leancharwidth}{M}\setlength{\parskip}{0pt}\setlength{\parindent}{0pt}
  \frenchspacing}{\par\endgroup\smallskip}
\newcommand{\leanline}[2]{\hangindent=\dimexpr#1\leancharwidth+4\leancharwidth\relax
  \hangafter=1\noindent\hspace*{#1\leancharwidth}#2\par}
\newcommand{\leansource}[1]{\par{\footnotesize\color{leanmuted}#1}\par}
\newcommand{\leanlabel}[1]{\par\smallskip\noindent{\small\bfseries #1}\par}
\usepackage[margin=25mm]{geometry}
\usepackage[unicode,hidelinks]{hyperref}
\hypersetup{pdfauthor={LangevinLean Project}}
\setlength{\emergencystretch}{3em}
\setlength{\parskip}{0.4em}
\setlength{\parindent}{0pt}
\allowdisplaybreaks[2]
\raggedbottom
\title{\texorpdfstring{Orthogonality of future Brownian increments to past observables}{Orthogonality of future Brownian increments to past observables}}
\author{LangevinLean Project}
\date{}
\begin{document}\maketitle
Draft mathematical exposition; these notes may contain mistakes.


Assume independent full scalar Brownian paths with measurable evaluations and initialization independent of the entire Brownian family. The integral solution is measurable with respect to initialization and the stopped noise path. A future Brownian increment is independent of this past and centered. Its expected inner product with any square-integrable Lipschitz observable of the reference state is therefore zero. This is an expectation identity, not a samplewise orthogonality claim.



Earlier results used here: \cite[Section~3.54]{LangevinLean2026ULASC}; \cite[Section~3.25]{LangevinLean2026ULASC}; \cite[Section~3.31]{LangevinLean2026ULASC}.

\section{\texorpdfstring{Independence of vector stopped past and shifted future}{Independence of vector stopped past and shifted future}}


For each scalar coordinate, the shifted future \(B_i(T+s)-B_i(T)\) is independent of the stopped past \(B_i(\min(t,T))\). Independence of the full coordinate paths gives joint independence across these coordinatewise pairs. Selecting and grouping the disjoint past and future indices therefore gives independence of the entire vector-valued past and future path families.



The stopped path records all coordinate values up to time \(t\); the shifted future records increments after \(t\). The resulting independence is between the sigma-algebras generated by these full path blocks. Consequently any measurable function of the entire past is independent of every measurable future increment. This uses full-path independence and Brownian independent increments, rather than unrelated pairwise independence assertions.

\begin{leanbox}{Lean: \texttt{brownian\_\allowbreak{}future\_\allowbreak{}past\_\allowbreak{}independent}}
\leansource{\texttt{brownian\_\allowbreak{}future\_\allowbreak{}past\_\allowbreak{}independent} \quad ProximalLangevin/\allowbreak{}SplittingOrthogonality.lean, lines 30–67}
\begin{leancode}
\leanline{0}{theorem\ brownian\_future\_past\_independent\ [Finite\ ι]}
\leanline{4}{(B\ :\ ι\ →\ ℝ≥0\ →\ Ω\ →\ ℝ)\ (hB\ :\ ∀\ i,\ IsPreBrownianReal\ (B\ i)\ Q)}
\leanline{4}{(hmeas\ :\ ∀\ i\ t,\ Measurable\ (B\ i\ t))}
\leanline{4}{(hpaths\ :\ iIndepFun\ (fun\ i\ ω\ t\ ↦\ B\ i\ t\ ω)\ Q)\ (T\ :\ ℝ≥0)\ :}
\leanline{4}{IndepFun\ (fun\ ω\ (i\ :\ ι)\ (s\ :\ ℝ≥0)\ ↦\ B\ i\ (T\ +\ s)\ ω\ -\ B\ i\ T\ ω)}
\leanline{6}{(fun\ ω\ (i\ :\ ι)\ (t\ :\ ℝ≥0)\ ↦\ B\ i\ (min\ t\ T)\ ω)\ Q}
\end{leancode}
\leanlabel{How the Lean proof follows the argument}
Formal correspondence: brownian\_future\_past\_independent : The complete shifted future and the stopped past of independent scalar Brownian coordinates are independent vector-valued path families.



\texttt{brownian\_\allowbreak{}future\_\allowbreak{}past\_\allowbreak{}independent}: the intermediate facts include \texttt{X}, \texttt{h\allowbreak{}Xm}, \texttt{hrow}, \texttt{hs}, \texttt{hr}, in source order; the source steps apply or rewrite \texttt{measurable\_\allowbreak{}pi\_\allowbreak{}iff.\allowbreak{}mpr}, \texttt{Bool.\allowbreak{}false\_\allowbreak{}eq\_\allowbreak{}true}, \texttt{hpaths.\allowbreak{}comp}, \texttt{Finset.\allowbreak{}disjoint\_\allowbreak{}left}, \texttt{hi.\allowbreak{}comp}.
\end{leanbox}

\section{\texorpdfstring{The pointwise Brownian vector and its square-integrable representative}{The pointwise Brownian vector and its square-integrable representative}}


Define the pointwise vector noise by \(\mathcal B_t=\sqrt2(B(t_+)-B(0))\) with \(t_+=\max(t,0)\). The constructed \(L^2\) Brownian noise agrees almost surely with this expression. Each scalar evaluation is measurable, and finite vector assembly, subtraction, and scaling preserve measurability. Its Gaussian coordinate second moments prove square integrability.



The \(L^2\) class and the pointwise vector thus have the same laws and integrals. Independence and centering may be transferred by their almost-sure equality. The nonnegative-time definition specifies the domain of Brownian evaluation and does not change the positive-time process.

\begin{leanbox}{Lean: \texttt{rawBrownianNoise}, \texttt{brownianNoise\_\allowbreak{}coe}, \texttt{rawBrownianNoise\_\allowbreak{}measurable}}
\leansource{\texttt{rawBrownianNoise} \quad ProximalLangevin/\allowbreak{}SplittingOrthogonality.lean, lines 70–71}
\begin{leancode}
\leanline{0}{def\ rawBrownianNoise\ (B\ :\ ι\ →\ ℝ≥0\ →\ Ω\ →\ ℝ)\ (t\ :\ ℝ)\ (ω\ :\ Ω)\ :\ EuclideanSpace\ ℝ\ ι\ :=}
\leanline{2}{Real.sqrt\ 2\ •\ (ULA.brownianCoordinates\ B\ t.toNNReal\ ω\ -\ ULA.brownianCoordinates\ B\ 0\ ω)}
\end{leancode}
\leansource{\texttt{brownianNoise\_\allowbreak{}coe} \quad ProximalLangevin/\allowbreak{}SplittingOrthogonality.lean, lines 74–85}
\begin{leancode}
\leanline{0}{theorem\ brownianNoise\_coe\ (B\ :\ ι\ →\ ℝ≥0\ →\ Ω\ →\ ℝ)}
\leanline{4}{(hB\ :\ ∀\ i,\ IsPreBrownianReal\ (B\ i)\ Q)\ (t\ :\ ℝ)\ :}
\leanline{4}{ULA.brownianLangevinNoise\ Q\ B\ hB\ 0\ t\ =ᵐ[Q]\ rawBrownianNoise\ B\ t}
\end{leancode}
\leansource{\texttt{rawBrownianNoise\_\allowbreak{}measurable} \quad ProximalLangevin/\allowbreak{}SplittingOrthogonality.lean, lines 88–94}
\begin{leancode}
\leanline{0}{theorem\ rawBrownianNoise\_measurable\ [Finite\ ι]\ (B\ :\ ι\ →\ ℝ≥0\ →\ Ω\ →\ ℝ)}
\leanline{4}{(hmeas\ :\ ∀\ i\ t,\ Measurable\ (B\ i\ t))\ (t\ :\ ℝ)\ :\ Measurable\ (rawBrownianNoise\ B\ t)}
\end{leancode}
\leanlabel{How the Lean proof follows the argument}
Formal correspondence: rawBrownianNoise : The pointwise vector presentation of the Langevin noise from time zero.



\texttt{raw\allowbreak{}Brownian\allowbreak{}Noise} is a definition with defining expression \texttt{Real.\allowbreak{}sqrt 2 \ensuremath{\bullet} (ULA.\allowbreak{}brownian\allowbreak{}Coordinates B t.\allowbreak{}to\allowbreak{}NNReal \ensuremath{\omega} -\allowbreak{} ULA.\allowbreak{}brownian\allowbreak{}Coordinates B 0 \ensuremath{\omega})}. The right-hand side specifies the object used in the mathematical statement.



\texttt{brownian\allowbreak{}Noise\_\allowbreak{}coe}: the source steps apply or rewrite \texttt{zero\_\allowbreak{}add}, \texttt{Real.\allowbreak{}to\allowbreak{}NNReal\_\allowbreak{}zero}, \texttt{Pi.\allowbreak{}smul\_\allowbreak{}apply}, \texttt{Pi.\allowbreak{}sub\_\allowbreak{}apply}.



\texttt{raw\allowbreak{}Brownian\allowbreak{}Noise\_\allowbreak{}measurable}: the source steps apply or rewrite \texttt{With\allowbreak{}Lp.\allowbreak{}measurable\_\allowbreak{}to\allowbreak{}Lp}, \texttt{measurable\_\allowbreak{}pi\_\allowbreak{}iff.\allowbreak{}mpr}.
\end{leanbox}


Earlier results used here: \cite[Section~3.24.14]{LangevinLean2026ULASC}; \cite[Section~3.24.15]{LangevinLean2026ULASC}; \cite[Section~3.24.21]{LangevinLean2026ULASC}; \cite[Section~3.24.8]{LangevinLean2026ULASC}; \cite[Section~3.25.1]{LangevinLean2026ULASC}.

\section{\texorpdfstring{Independence of the initialized past from future increments}{Independence of the initialized past from future increments}}


The stopped Brownian path is independent of the future increment. Initialization is independent of their joint pair. Reassociation of the product law therefore makes the pair consisting of initialization and stopped past independent of the future increment. Any variable measurable with respect to that initialized past inherits the same independence.



A variable only almost everywhere strongly measurable for the past admits a measurable representative with the same law. Almost-sure equality then transfers the conclusion to its \(L^2\) representative. The grouping relation is \((\text{initial},\text{past})\mid\text{future}\). Thus independence for the later reference observable follows from its past measurability and the input assumptions; it is not added as an independent assertion about the solution.

\begin{leanbox}{Lean: \texttt{brownian\_\allowbreak{}future\_\allowbreak{}independent\_\allowbreak{}initial\_\allowbreak{}past}, \texttt{naturalPast\_\allowbreak{}independent\_\allowbreak{}of\_\allowbreak{}aestronglyMeasurable}}
\leansource{\texttt{brownian\_\allowbreak{}future\_\allowbreak{}independent\_\allowbreak{}initial\_\allowbreak{}past} \quad ProximalLangevin/\allowbreak{}SplittingOrthogonality.lean, lines 98–171}
\begin{leancode}
\leanline{0}{theorem\ brownian\_future\_independent\_initial\_past}
\leanline{4}{(B\ :\ ι\ →\ ℝ≥0\ →\ Ω\ →\ ℝ)\ (hB\ :\ ∀\ i,\ IsPreBrownianReal\ (B\ i)\ Q)}
\leanline{4}{(hmeas\ :\ ∀\ i\ t,\ Measurable\ (B\ i\ t))}
\leanline{4}{(hpaths\ :\ iIndepFun\ (fun\ i\ ω\ t\ ↦\ B\ i\ t\ ω)\ Q)}
\leanline{4}{(initial\ :\ Lp\ (EuclideanSpace\ ℝ\ ι)\ 2\ Q)}
\leanline{4}{(hinit\ :\ IndepFun\ initial\ (fun\ ω\ (i\ :\ ι)\ (t\ :\ ℝ≥0)\ ↦\ B\ i\ t\ ω)\ Q)}
\leanline{4}{(t\ h\ :\ ℝ)\ (ht\ :\ 0\ ≤\ t)\ (hth\ :\ t\ ≤\ h)\ :}
\leanline{4}{IndepFun}
\leanline{6}{(fun\ ω\ (r\ :\ Option\ (Icc\ (0\ :\ ℝ)\ t))\ ↦}
\leanline{8}{r.elim\ (initial\ ω)\ (fun\ s\ ↦\ ULA.brownianLangevinNoise\ Q\ B\ hB\ 0\ s\ ω))}
\leanline{6}{(ULA.brownianLangevinNoise\ Q\ B\ hB\ 0\ h\ -\ ULA.brownianLangevinNoise\ Q\ B\ hB\ 0\ t)\ Q}
\end{leancode}
\leansource{\texttt{naturalPast\_\allowbreak{}independent\_\allowbreak{}of\_\allowbreak{}aestronglyMeasurable} \quad ProximalLangevin/\allowbreak{}SplittingOrthogonality.lean, lines 176–196}
\begin{leancode}
\leanline{0}{theorem\ naturalPast\_independent\_of\_aestronglyMeasurable}
\leanline{4}{(initial\ :\ Lp\ (EuclideanSpace\ ℝ\ ι)\ 2\ Q)}
\leanline{4}{(noise\ :\ ℝ\ →\ Lp\ (EuclideanSpace\ ℝ\ ι)\ 2\ Q)\ (t\ :\ ℝ)}
\leanline{4}{(W\ :\ Ω\ →\ EuclideanSpace\ ℝ\ ι)}
\leanline{4}{(hind\ :\ IndepFun\ (fun\ ω\ (r\ :\ Option\ (Icc\ (0\ :\ ℝ)\ t))\ ↦}
\leanline{6}{r.elim\ (initial\ ω)\ (fun\ s\ ↦\ noise\ s\ ω))\ W\ Q)}
\leanline{4}{(Z\ :\ Ω\ →\ EuclideanSpace\ ℝ\ ι)}
\leanline{4}{(hZ\ :\ AEStronglyMeasurable[ULA.lpLangevinPast\ Q\ initial\ noise\ t]\ Z\ Q)\ :}
\leanline{4}{IndepFun\ Z\ W\ Q}
\end{leancode}
\leanlabel{How the Lean proof follows the argument}
Formal correspondence: brownian\_future\_independent\_initial\_past : A future Brownian increment is independent of the initialization and all earlier actual \(L^2\) Brownian-noise representatives. naturalPast\_independent\_of\_aestronglyMeasurable : Every random variable measurable from the natural initialization/\allowbreak{}noise past inherits its independence from a supplied future source variable.



\texttt{brownian\_\allowbreak{}future\_\allowbreak{}independent\_\allowbreak{}initial\_\allowbreak{}past}: the intermediate facts include \texttt{hhh}, \texttt{hdelta}, \texttt{hadd}, \texttt{hm}, \texttt{hopt}, in source order; the source steps apply or rewrite \texttt{NNReal.\allowbreak{}eq}, \texttt{NNReal.\allowbreak{}coe\_\allowbreak{}add}, \texttt{brownian\allowbreak{}Noise\_\allowbreak{}coe}, \texttt{hopt'.\allowbreak{}congr}, \texttt{hd.\allowbreak{}trans}. The remaining scalar rearrangements are checked by \texttt{ring}.



\texttt{natural\allowbreak{}Past\_\allowbreak{}independent\_\allowbreak{}of\_\allowbreak{}aestrongly\allowbreak{}Measurable}: the intermediate facts include \texttt{F}, \texttt{hs}, \texttt{hm}, in source order; the source steps apply or rewrite \texttt{sup\_\allowbreak{}le}, \texttt{Measurable\allowbreak{}Space.\allowbreak{}comap\_\allowbreak{}le\_\allowbreak{}comap\_\allowbreak{}pi}, \texttt{hm.\allowbreak{}congr}, \texttt{h\allowbreak{}Z.\allowbreak{}ae\_\allowbreak{}eq\_\allowbreak{}mk.\allowbreak{}symm}, \texttt{Filter.\allowbreak{}Eventually\allowbreak{}Eq.\allowbreak{}rfl}.
\end{leanbox}


Earlier results used here: \cite[Section~3.24.8]{LangevinLean2026ULASC}; \cite[Section~3.25.1]{LangevinLean2026ULASC}; \cite[Section~3.54.4]{LangevinLean2026ULASC}.

\section{\texorpdfstring{Expected orthogonality from independence and centering}{Expected orthogonality from independence and centering}}


Scalar Brownian centering and commutation of finite coordinate projections with integration give \(\mathbb E B_t=0\), and subtraction and scaling give \(\mathbb E\mathcal B_t=0\). For independent \(L^2\) vectors \(X,W\), Cauchy–Schwarz proves integrability of the inner product. Independence and bilinear integration give \(\langle X,W\rangle_{L^2}=\langle\mathbb E X,\mathbb E W\rangle\).



When \(W\) is centered, the latter inner product of means is zero. Thus no norm bound on the other mean is needed: independence and centering give exact expected orthogonality.

\begin{leanbox}{Lean: \texttt{brownianCoordinates\_\allowbreak{}integral\_\allowbreak{}zero}, \texttt{brownianNoise\_\allowbreak{}integral\_\allowbreak{}zero}, \texttt{independent\_\allowbreak{}centered\_\allowbreak{}lp\_\allowbreak{}inner\_\allowbreak{}zero}}
\leansource{\texttt{brownianCoordinates\_\allowbreak{}integral\_\allowbreak{}zero} \quad ProximalLangevin/\allowbreak{}SplittingOrthogonality.lean, lines 199–209}
\begin{leancode}
\leanline{0}{theorem\ brownianCoordinates\_integral\_zero}
\leanline{4}{(B\ :\ ι\ →\ ℝ≥0\ →\ Ω\ →\ ℝ)\ (hB\ :\ ∀\ i,\ IsPreBrownianReal\ (B\ i)\ Q)\ (t\ :\ ℝ≥0)\ :}
\leanline{4}{(∫\ ω,\ ULA.brownianCoordinates\ B\ t\ ω\ ∂Q)\ =\ 0}
\end{leancode}
\leansource{\texttt{brownianNoise\_\allowbreak{}integral\_\allowbreak{}zero} \quad ProximalLangevin/\allowbreak{}SplittingOrthogonality.lean, lines 212–223}
\begin{leancode}
\leanline{0}{theorem\ brownianNoise\_integral\_zero}
\leanline{4}{(B\ :\ ι\ →\ ℝ≥0\ →\ Ω\ →\ ℝ)\ (hB\ :\ ∀\ i,\ IsPreBrownianReal\ (B\ i)\ Q)\ (t\ :\ ℝ)\ :}
\leanline{4}{(∫\ ω,\ ULA.brownianLangevinNoise\ Q\ B\ hB\ 0\ t\ ω\ ∂Q)\ =\ 0}
\end{leancode}
\leansource{\texttt{independent\_\allowbreak{}centered\_\allowbreak{}lp\_\allowbreak{}inner\_\allowbreak{}zero} \quad ProximalLangevin/\allowbreak{}SplittingOrthogonality.lean, lines 226–236}
\begin{leancode}
\leanline{0}{theorem\ independent\_centered\_lp\_inner\_zero}
\leanline{4}{(X\ W\ :\ Lp\ (EuclideanSpace\ ℝ\ ι)\ 2\ Q)\ (hind\ :\ IndepFun\ X\ W\ Q)}
\leanline{4}{(hmean\ :\ (∫\ ω,\ W\ ω\ ∂Q)\ =\ 0)\ :\ inner\ ℝ\ X\ W\ =\ 0}
\end{leancode}
\leanlabel{How the Lean proof follows the argument}
Formal correspondence: brownianCoordinates\_integral\_zero : The finite-dimensional Brownian vector has zero mean. brownianNoise\_integral\_zero : The actual \(L^2\) additive Langevin noise is centered at every real time. independent\_centered\_lp\_inner\_zero : Binary independence and zero mean imply genuine \(L^2\) orthogonality.



\texttt{brownian\allowbreak{}Coordinates\_\allowbreak{}integral\_\allowbreak{}zero}: the intermediate facts include \texttt{hint}, \texttt{hi}, in source order; the source steps apply or rewrite \texttt{integral\_\allowbreak{}eval}.



\texttt{brownian\allowbreak{}Noise\_\allowbreak{}integral\_\allowbreak{}zero}: the intermediate fact is \texttt{hint}, in source order; the source steps apply or rewrite \texttt{integral\_\allowbreak{}congr\_\allowbreak{}ae}, \texttt{brownian\allowbreak{}Noise\_\allowbreak{}coe}, \texttt{integral\_\allowbreak{}smul}, \texttt{integral\_\allowbreak{}sub}.



\texttt{independent\_\allowbreak{}centered\_\allowbreak{}lp\_\allowbreak{}inner\_\allowbreak{}zero}: the intermediate fact is \texttt{hi}, in source order; the source steps apply or rewrite \texttt{L2.\allowbreak{}inner\_\allowbreak{}def}, \texttt{inner\_\allowbreak{}zero\_\allowbreak{}right}.
\end{leanbox}


Earlier results used here: \cite[Section~3.24.1]{LangevinLean2026ULASC}; \cite[Section~3.24.8]{LangevinLean2026ULASC}; \cite[Section~3.24.9]{LangevinLean2026ULASC}; \cite[Section~3.25.1]{LangevinLean2026ULASC}.

\section{\texorpdfstring{Orthogonality for observables of the adapted integral solution}{Orthogonality for observables of the adapted integral solution}}


The integral-equation construction makes \(Y_t\) measurable for the initialized Brownian past. A Lipschitz observable \(b(Y_t)\) remains measurable for that past and belongs to \(L^2\). Independence of future increments and their zero mean imply

\[\langle b(Y_t),\mathcal B_h-\mathcal B_t\rangle_{L^2}=0,\qquad0\le t\le h.\]

For observable \(b(Y_t)\), both factors are in \(L^2\), so their pairing is integrable and

\[\mathbb E\langle b(Y_t),B_h-B_t\rangle=\left\langle\mathbb E b(Y_t),\mathbb E(B_h-B_t)\right\rangle=0.\]

The \(L^2\) inner product is this expectation. The state at time \(t\) depends only on initialization and noise up to \(t\). Composing with the drift observable preserves this dependence, while the increment over \([t,h]\) is independent and centered. This identity removes the contribution that an absolute norm bound would estimate at order \(\sqrt{h-t}\) in the subsequent energy comparison.

\begin{leanbox}{Lean: \texttt{lp\_\allowbreak{}langevin\_\allowbreak{}future\_\allowbreak{}inner\_\allowbreak{}zero}}
\leansource{\texttt{lp\_\allowbreak{}langevin\_\allowbreak{}future\_\allowbreak{}inner\_\allowbreak{}zero} \quad ProximalLangevin/\allowbreak{}SplittingOrthogonality.lean, lines 241–277}
\begin{leancode}
\leanline{0}{theorem\ lp\_langevin\_future\_inner\_zero}
\leanline{4}{(U\ :\ EuclideanSpace\ ℝ\ ι\ →\ ℝ)\ (L\ :\ ℝ≥0)\ (hLip\ :\ LipschitzWith\ L\ (gradient\ U))}
\leanline{4}{(initial\ :\ Lp\ (EuclideanSpace\ ℝ\ ι)\ 2\ Q)}
\leanline{4}{(B\ :\ ι\ →\ ℝ≥0\ →\ Ω\ →\ ℝ)\ (hB\ :\ ∀\ i,\ IsPreBrownianReal\ (B\ i)\ Q)}
\leanline{4}{(hmeas\ :\ ∀\ i\ s,\ Measurable\ (B\ i\ s))}
\leanline{4}{(hpaths\ :\ iIndepFun\ (fun\ i\ ω\ s\ ↦\ B\ i\ s\ ω)\ Q)}
\leanline{4}{(hinit\ :\ IndepFun\ initial\ (fun\ ω\ (i\ :\ ι)\ (s\ :\ ℝ≥0)\ ↦\ B\ i\ s\ ω)\ Q)}
\leanline{4}{(Y\ :\ ℝ\ →\ Lp\ (EuclideanSpace\ ℝ\ ι)\ 2\ Q)\ (hY\ :\ Continuous\ Y)}
\leanline{4}{(t\ h\ :\ ℝ)\ (ht\ :\ 0\ ≤\ t)\ (hth\ :\ t\ ≤\ h)}
\leanline{4}{(hYe\ :\ ∀\ s\ ∈\ Icc\ (0\ :\ ℝ)\ h,\ Y\ s\ -\ initial\ =}
\leanline{6}{-(∫\ r\ in\ 0..s,\ ULA.lpDrift\ Q\ (gradient\ U)\ hLip\ (Y\ r))\ +}
\leanline{8}{ULA.brownianLangevinNoise\ Q\ B\ hB\ 0\ s)}
\leanline{4}{(b\ :\ EuclideanSpace\ ℝ\ ι\ →\ EuclideanSpace\ ℝ\ ι)\ (Lb\ :\ ℝ≥0)}
\leanline{4}{(hLipb\ :\ LipschitzWith\ Lb\ b)\ :}
\leanline{4}{inner\ ℝ\ (ULA.lpDrift\ Q\ b\ hLipb\ (Y\ t))}
\leanline{6}{(ULA.brownianLangevinNoise\ Q\ B\ hB\ 0\ h\ -\ ULA.brownianLangevinNoise\ Q\ B\ hB\ 0\ t)\ =\ 0}
\end{leancode}
\leanlabel{How the Lean proof follows the argument}
Formal correspondence: lp\_langevin\_future\_inner\_zero : A Lipschitz function of the actual Langevin solution is orthogonal to every later Brownian increment. Adaptedness is derived from the integral equation by the existing invariant-subspace construction.



\texttt{lp\_\allowbreak{}langevin\_\allowbreak{}future\_\allowbreak{}inner\_\allowbreak{}zero}: the intermediate facts include \texttt{noise}, \texttt{hadapt}, \texttt{h\allowbreak{}Z}, \texttt{hi\ensuremath{{}_{1}}}, \texttt{hi\ensuremath{{}_{2}}}, in source order; the source steps apply or rewrite \texttt{independent\_\allowbreak{}centered\_\allowbreak{}lp\_\allowbreak{}inner\_\allowbreak{}zero}, \texttt{integral\_\allowbreak{}congr\_\allowbreak{}ae}, \texttt{integral\_\allowbreak{}sub}, \texttt{brownian\allowbreak{}Noise\_\allowbreak{}integral\_\allowbreak{}zero}, \texttt{sub\_\allowbreak{}self}.
\end{leanbox}


Earlier results used here: \cite[Section~3.13.2]{LangevinLean2026ULASC}; \cite[Section~3.25.1]{LangevinLean2026ULASC}; \cite[Section~3.25.2]{LangevinLean2026ULASC}; \cite[Section~3.25.3]{LangevinLean2026ULASC}; \cite[Section~3.54.2]{LangevinLean2026ULASC}; \cite[Section~3.54.4]{LangevinLean2026ULASC}; \cite[Section~3.54.8]{LangevinLean2026ULASC}.

\begin{thebibliography}{ULA-LSI}
\bibitem[ULA-SC]{LangevinLean2026ULASC}
LangevinLean Project. \emph{Unadjusted Langevin under Strong Convexity}. LL-01, 2026. Draft report, content revision d34ead06987c. \url{https://odakyildiz.github.io/LangevinLean/docs/pdf/UnadjustedLangevin/FullProof.pdf}.
\end{thebibliography}

\end{document}
